%!TEX root = ../thesis.tex
% Chapter: Discussion

\chapter{Discussion} \label{sec:discussion}

The paper-style replication reproduces all three of Razali et al.'s configurations on
both tasks (Chapter~\ref{sec:results}). The project's main result lies beyond it,
in the mass extensions (Chapter~\ref{sec:extensions}): preserving the
wavelet-scattering spatial map is the decisive lever. This chapter interprets the replication against the paper, draws
out what the extensions establish, critiques the paper's single-split evaluation protocol
and defines the leakage-free one used here, and sets out the data-provenance caveats and
limitations that qualify the numbers.

\section{Interpreting the replication: where fusion helps, and why}

The CNN~+~WS fusion evaluated here is a faithful instance of the paper's Model~5: a
training-fold-standardised concatenation of the $512$-dimensional CNN average-pool vector
and the $406$-dimensional wavelet-scattering vector, classified by the same Subspace k-NN
ensemble, with no GLCM branch. Read against the four implementation differences in Chapter~\ref{sec:results}, the
interpretation is clear: the paper's central mechanism (handcrafted scattering features
rescuing a data-starved CNN) reproduces \emph{exactly in the regime the paper invokes it,
and only there}. On
\textbf{tissue}, our CNN does not overfit (test ${\sim}98\%$, AU-ROC $0.979$; an overfit
gap of ${\sim}2$~pp against the paper's $21.6$~pp), so fusion has almost nothing to add:
CNN~+~WS ($0.988$) is statistically close to CNN-only ($0.979$), and the large fusion
gain the paper reports ($77.4\% \rightarrow 93.5\%$ test) does not appear, because the
deficit it compensates is absent under our cleaner preprocessing. On \textbf{mass}, the
CNN genuinely overfits (training accuracy saturates by epoch~10 while validation
oscillates in the $60$--$90\%$ band, Figure~\ref{fig:mass_overfit}) and adding the
training-free scattering features lifts the leakage-free result: fusion (AU-ROC
$0.816$) outranks both CNN-only ($0.762$) and WS-only ($0.726$). This is where fusion helps: it reproduces the paper's qualitative claim
on the mass task, and shows precisely why that claim does not transfer to our
tissue pipeline.

The paper's headline mass fusion result rests on a single ${\sim}22$-patch split; our
cross-validation reproduces its ${\approx}98\%$, but the repeated held-out estimate is
lower and noisier ($0.816 \pm 0.057$). Why that $1.00$ is not a stable target is taken
up in Section~\ref{sec:eval_disc}.

\section{The WS-only tissue gap and what closes it}

The single largest replication discrepancy is the ${\sim}7$~pp shortfall on WS-only tissue
(AU-ROC $0.895$ against the paper's $0.94$). Of the plausible causes, density-aware
contrast enhancement, the Kymatio-versus-MATLAB library convention ($406$ versus $391$
coefficients), and ACR-guided patch sampling, we cannot fully disentangle the three, but
contrast enhancement is the leading candidate: wavelet scattering is a textural descriptor of
frequency content, and the paper's density-aware histogram adjustment amplifies exactly the
density-specific contrast the tissue task measures, whereas our density-agnostic CLAHE does not. The more consequential
observation, however, is that this gap is closed and substantially exceeded by
changing the representation rather than the preprocessing: on the same patches and pipeline,
the adapted scattering covariance (Section~\ref{sec:scov_tissue}) reaches AU-ROC $0.984$,
$+0.089$ over averaged WS, because the standard $S_0, S_1, S_2$ coefficients undersample the
cross-scale and cross-orientation structure that covariance exposes. The ceiling on the WS
branch for this project is therefore set by the richness of the descriptor, not by the
contrast operation. On mass, by contrast, the same richer scalar descriptor gives no
gain ($0.734$ versus $0.726$), foreshadowing the central extension result below.

\section{What the extensions establish}

The extensions establish one thing for the mass task: \emph{preserving
spatial and lesion structure in the wavelet-scattering representation matters more than
changing the scalar features or the classifier}. The screen of
Section~\ref{sec:screening} shows that normalisation, per-path statistics, coarse pooling,
feature selection, and even the richer scattering-covariance descriptor move mass AU-ROC
only within across-seed noise of the $0.726$ averaged-WS baseline; the one method that
breaks out is the full-patch WS$\rightarrow$CNN head ($0.909$), which keeps the $J=3$
spatial maps and learns over them instead of averaging them away. Conditioning those maps
on lesion morphology lifts the classical model from $0.801$ to $0.892$, and a learned CNN
head over the region-aware maps reaches $0.946$. A lighter, masked/no-CLAHE input lifts this
region-aware head further to AU-ROC $0.971$ (Section~\ref{sec:region_noclahe}), consistent with
CLAHE's local contrast enhancement distorting the scattering energy the descriptor relies on.

The contrast with tissue explains why. Tissue density is a texture problem that
a richer scalar descriptor already solves ($0.984$), with no need to retain spatial
layout; mass is a structured-object problem in which \emph{where} the signal sits (core
versus margin versus context) is decisive, so the scalar route plateaus and the spatial
route succeeds. Region-aware preservation is therefore the project's principal extension on
the mass side. The CNN and fusion methods reinforce that
conclusion rather than overturn it: the region-aware late fusion is numerically strongest
($0.948$), but it is a decision-level combination of two trained models whose margin over
the region-aware WS$\rightarrow$CNN alone is within seed noise, so it is reported as the
best \emph{number}, while the region-aware representation is the real contribution.

\section{Evaluation: the paper's protocol and ours} \label{sec:eval_disc}

\textbf{The paper's point-estimate protocol.} A methodological weakness of the paper
compounds the small mass sample: it reports single-split point estimates, most
strikingly a mass CNN~+~WS AU-ROC of $1.00$ and $91.7\%$ accuracy, on a held-out set of
only ${\sim}22$ patches, with no measure of variance. At that size a single misclassification
shifts accuracy by ${\sim}4$~pp, and across our five seeds the same method's held-out mass
AU-ROC ranges from roughly $0.70$ to $0.76$; a lone favourable split can therefore read as
$1.00$. A point estimate at this sample size is neither reproducible nor a meaningful
measure of generalisation. This is the central reason the project does not report
single-split numbers, and it motivates the evaluation protocol used throughout
(Section~\ref{sec:eval_protocol}).

\textbf{Repeated, leakage-free splits.} Every headline mass method is reported as a
five-seed mean~$\pm$~standard deviation over image-level splits; the residual bands
(e.g.\ $\pm0.036$ for the region-aware WS$\rightarrow$CNN, and WS-only mass accuracy ranging
from $0.64$ to $0.83$ across the five seeds against a mean of $0.701$) reflect the small sample honestly rather than
concealing it. Splits are made at the image level with group-aware cross-validation, so
patches from one mammogram never straddle the train/test boundary; the paper does not
specify its split, so we adopt the leakage-free convention throughout
(Section~\ref{sec:stage3}).

\textbf{Determinism of learned models.} Python, NumPy, PyTorch, and cuDNN seeds and flags
are pinned for the learned models, but the WS-map CNN is still trained on only ${\sim}90$
patches; the mitigation is to average five seeds and corroborate with the tighter
group-aware cross-validation ($0.919 \pm 0.012$), which is far less split-sensitive than any
single held-out split.

\textbf{Optimistic CV on learned features.} Cross-validation computed on learned CNN
or fusion features is optimistic, because the extractor has effectively already seen the
inner-validation fold; the ${\sim}0.99$ mass CV figures (marked ``(opt.)'' in
Table~\ref{tab:results_mass}) are not out-of-sample estimates, whereas the five-seed
held-out AU-ROCs ($0.762$ and $0.816$) are. The project therefore leads with held-out AU-ROC
as the mass headline and uses group-aware CV only as a stability check, and, for the
classical WS methods where nothing is learned, as a like-for-like comparison with the paper.

\section{Data provenance and quality}

Three data-provenance issues surfaced during the project, and none changes the qualitative conclusions.

\textbf{Mass patch count.} \label{sec:mass_count_discussion} Our reproducible \texttt{Name="Mass"} XML extraction yields $116$
patches across $107$ images ($75$ malignant, $41$ benign), against the paper's $112$
($76+36$). The difference arises from two editorial choices the paper leaves undocumented:
we retain all mass ROIs (eight images carry more than one: a mix of near-duplicate
annotations and genuinely multifocal disease), and we admit only ROIs literally tagged
\texttt{Mass}, whereas the paper appears to fold in a few architectural-distortion and
asymmetry findings. Ours is exactly reproducible from the public XML; the paper's requires
per-image judgement. Because the split is image-level, the extra ROIs add slight within-fold
redundancy but no leakage, and the difference sits well below the mass test-set noise floor.

\textbf{Data-quality filter.} \label{sec:mass_quality_filter} One annotation,
\texttt{file\_id 50994354}, BI-RADS~4a, view \texttt{MG\_R\_ML}, has a bounding box that
falls entirely outside the breast mask (an annotation/orientation mismatch unique to this
file), so cropping it from the masked image gives a uniformly-zero patch. A dynamic one-line
filter (\texttt{patch.max() > 0}) excludes it, reducing the mass set to the \textbf{115
patches} ($74$ malignant, $41$ benign) the replication tables use (Chapter~\ref{sec:results});
the filter is general, so it would catch any future analogous patch. The on-disk patches and
the upstream index are unchanged.

\textbf{Annotation-source disagreements.} \label{sec:annot_audit} INbreast ships two
independent versions of each annotation the pipeline consumes, and they do not perfectly
agree. For lesion presence, the XLS \texttt{Mass} column flags $108$ files and the XML
\texttt{Mass} ROIs $107$, agreeing on $106$; the three disagreements are two low-acuity
BI-RADS-2 XLS-only cases and one high-acuity BI-RADS-4c XML-only case (\texttt{24065584}).
The pipeline takes ROIs from the XML because it needs polygon coordinates, which happens to
trade two benign cases for one high-acuity malignant one, a shift well below the mass
test-set noise floor. For pectoral masks, the
LabelMe-JSON source ($202$ MLO views) and the Apple-plist XML source ($201$) agree on $199$
files and differ only in outline style; the choice has negligible downstream effect because
the Otsu breast mask and the patch sampler's minimum target-fraction absorb the difference.

\section{Limitations}

\begin{itemize}
    \item \textbf{Small mass sample.} The mass dataset is only $115$ patches (${\sim}24$ held
    out per split), so the mass AU-ROC bands stay wide even averaged over five seeds; this is
    the dominant limitation on the mass side, and any clinical claim about mass malignancy
    would require external, multi-site validation before it could be trusted.
    \item \textbf{Single dataset.} Only INbreast is used; cross-dataset generalisation (e.g.\
    DDSM, MIAS) is untested.
    \item \textbf{No augmentation.} Following the paper, no augmentation is used: the most
    obvious lever left for the observed CNN overfitting. Because augmentation could have narrowed
    the very mass-CNN overfitting that fusion compensates, part of the measured fusion benefit
    may reflect this choice rather than an intrinsic advantage of handcrafted features.
    \item \textbf{Feature standardisation.} The paper does not mention standardising features
    before the Subspace k-NN; we do, since k-NN is distance-based. This may slightly help our
    WS results, but the qualitative picture (WS-tissue trailing the paper) is unchanged.
    \item \textbf{GLCM branch excluded.} The GLCM and CNN+GLCM/CNN+WS+GLCM configurations
    (paper Models~3, 4, 6) are out of scope; the paper itself reports GLCM adds no
    accuracy, precision, recall or F1 and only $+0.01$ AU-ROC over CNN+WS, so the practical
    impact is small.
    \item \textbf{Mass invariance scale not fully tuned.} The scattering-covariance $J/L$
    tuning was carried out on tissue; on mass we ran only an $L$-control at $J=6$ (no gain)
    and adopted $J=3$ maps, consistent with Bruna and Mallat's \cite{bruna2013invariant}
    observation that object-like tasks favour a smaller invariance scale that retains spatial
    structure. A full mass-specific $J\times L$ grid was not run, but the region-aware results
    indicate the larger lever for mass is spatial structure, not the scalar invariance scale.
\end{itemize}
